\documentclass[a4paper,12pt]{article}
\usepackage[utf8]{inputenc}
\usepackage[russian]{babel}
\usepackage{amsmath, amssymb}
\usepackage[hyperfootnotes=false]{hyperref}
\usepackage{geometry}
\usepackage{indentfirst}
\usepackage{graphicx}
\usepackage{float}
\usepackage{enumitem}
\usepackage{wrapfig}
\usepackage{caption}
\usepackage{array}
\geometry{
	a4paper,
	top=24mm,
	bottom=30mm,
	left=20mm,
	right=20mm
}

\newcommand{\blfootnote}[1]{%
	\begingroup
	\renewcommand{\thefootnote}{}%
	\footnote{#1}%
	\addtocounter{footnote}{-1}%
	\endgroup
}

% Отступ первой строки абзаца 1 см
\setlength{\parindent}{1cm}
\setlength{\parskip}{0pt}
\linespread{1}                 % одинарный интервал

% Отключение нумерации страниц
\pagenumbering{gobble}

% Настройка подписей рисунков и таблиц
% Используем \caption* для отключения точки, но оставляем обычный \caption
\usepackage[labelsep=space, justification=centering]{caption}
\captionsetup[figure]{font=small}   % small = 11pt при 12pt основного шрифта
\captionsetup[table]{font=small}
\renewcommand{\figurename}{Рис.}
\renewcommand{\tablename}{Табл.}

% Убираем точку в конце подписи (переопределяем \caption)
\usepackage{etoolbox}
\makeatletter
\renewcommand{\fnum@figure}{\figurename~\thefigure}
\renewcommand{\fnum@table}{\tablename~\thetable}
\apptocmd{\@caption}{\def\@captionheadfont{\normalfont}\def\@captionfont{\normalfont}}{}{}
\makeatother

% Настройка списка литературы
\makeatletter
\renewcommand{\@biblabel}[1]{#1.}   % нумерация "1.", "2."
\makeatother

\begin{document}
	\begin{flushright}
		И. М. Куликовских$^{1}$\blfootnote{*) И. М. Куликовских, kulikovskikh@bk.ru}
	\end{flushright}
	
	
	\begin{flushright}
		$^{1}$Санкт-Петербургский политехнический университет Петра Великого
	\end{flushright}

	\begin{center}
		ПЛОЩАДИ И ЛОГРАФИМЫ
	\end{center}
	
	\quad \textbf{Введение.} Лекция "Площади и логарифмы"  была прочитана Алексеем Ивановичем Маркушевичем осенью 1951 года в Московском университете для школьников 9 и 10 классов. Цель лекции -- изложить геометрическую теорию логарифмов к всесоюзной олимпиаде по математике. К сожалению, в современной старшей школе программа по математике, зачастую, подаётся кусочно и обрывисто, поэтому в данной работе хочется рассказать читателям о таком красивом и необычном подходе к логарифмам. Эта работа будет полезна к ознакомлению как учителям так и ученикам старшей школы, здесь работа А. И. Маркушевича будет представлена переосмысленной на новые образовательные тенденции. 
	
	\quad \textbf{\S 1. Площадь криволинейной трапеции.} 
	
	\quad Задача о нахождении площади криволинейной трапеции не представляет сложности в том случае если читатель знаком с интегральным счислением и понимает что такое определённый интеграл, однако существуют красивые частные случаи в которых совершенно не обязательно уметь интегрировать.
	
	\quad Рассмотрим степенную функицю:
	
	\begin{equation}
		f(x) = x^n \quad \forall n \in \mathbb{N}
	\end{equation}
	
	
	\quad Данное семейство функуций обладладает рядом полезных свойств:
	
	\begin{enumerate}
		\item Степенная функция с натуральным показателем непрерывна на всей вещественной оси.
		
		\item Монотонность зависит от чётности $n$, иначе говоря функция либо имеет постоянный характер монотонности всюду на $\mathbb{R}$, либо характер монотонности меняется лишь один раз.
	\end{enumerate}
	
	\begin{wrapfigure}[15]{r}{0.4\textwidth}
		\includegraphics[width=0.4\textwidth]{1.png}
		\caption{Криволинейная трапеция}
	\end{wrapfigure}
	
	\quad \underline{Опр.} Зафиксируем $n=2$. Этот случай называется параболой. Будем рассматривать лишь тот участок параболы, где функция не меняет характер монотонности (возрастания/убывания). Выделим отрезок $[a, b]$ (рис. 1), тогда фигура $ABCD$ называется \textit{криволинейной трапецией}. О нахождении площади этой фигуры говорят как о нахождении площади под графиком функции $f(x)$. 
	
	
	\quad Поставим задачу о вычислении площади криволинейной трапеции для всякой степенной функции вида (1). Так как в рамках школьной программы нет никаких инструментов решения подобной задачи, то предлагается искать площадь оценочно. Покажем что при достаточно сильном приближении ошибка между оценкой и реальным значением площади стремится к нулю. 
	
	\begin{wrapfigure}[15]{l}{0.4\textwidth}
		\includegraphics[width=0.4\textwidth]{2.png}
		\caption{Разбиение площади}
	\end{wrapfigure}
	
	\quad Разобьём площадь под графиком функции $f(x)$ прямоугольниками так, как показано на рисунке 2. В этом случае будем считать, что наша оценка площади равна сумме площадей всех прямоугольников. Ошибка между оценкой и реальным значением площади равна сумме площадей всех криволинейных треугольников. Возьмём прямоугольник с самой большой площадью из всех в разбиении. Окажется что все криволинейные треугольники полностью вмещаются в этот прямоугольник. Увеличивая число прямоугольников в разбиении при этом уменьшая площадь самого большого прямоугольника мы получаем сколь угодно малую ошибку между площадью и оценкой площади. Пусть это рассуждение пока не даёт конктреных аналитических выражений для площади криволинейной трапеции, зато позволяет провести следующее замечение:
	
	\quad \underline{Зам.:} Такое построение не зависит от вида разбиения. Каково бы ни было изначальное соотношение между площадями прямоугольников, измельчая дробление вышеуказанными способом мы уточняем оценку. 

	\quad Теперь мы можем переходить непосредственно к поиску площади криволинейной трапеции. Пользуясь фактом о том, что нам известен алгоритм уточнения оценки площади, мы можем подобрать конкретный частный случай с известной начальной оценкой и заданым алгоритмом уточнения. Введём некоторые определения:

	\quad \underline{Опр.} Пусть $\tau = \{x_k\}_{k=1}^{n}$ разбиение отрезка $[a, b]$ т. е. это множество точек $\{x_{k}\}_{k=1}^{n}$ такое что $a = x_1 < x_2 < x_3 < ... < x_n = b$, $\Delta x_k = x_k - x_{k-1}$, $|\tau| = max\{\Delta x_k\}$ -- мелкость разбиения $\tau$. 

	\quad \underline{Опр.} $S_{[a, b]} (f(x)) $ есть площадь под графиком функции $f(x)$ вида (1) на отрезке $[a, b]$. 
	
	\quad \underline{Теорема:} Рассмотрим луч $x\in(0, +\infty)$ и заданную на нём функцию $f(x)$ вида (1), тогда: 
	
	\begin{equation}
		S_{[a, b]} (f(x)) = \dfrac{b^{k+1} - a^{k+1}}{k+1}
	\end{equation}
	
	
	\quad \underline{Доказательство. }

	
	
	зам. : можно разбить на участки по монотонности и затем применить формулу для каждого из участков. 
	
	\quad \textbf{\S 2. Степенная функция с целым показателем. }
	
	\quad Для нахождения площади криволинейной трапеции всякой степенной функции с целым показателем достаточно определиться с нахождением площади под 
	
	теорема логарифм есть площадь под графиком (определена как спец функция)
	
	\quad \textbf{\S 3. Свойства натурального логарифма}
	
	\quad \textbf{Заключение}
	
\end{document}